\documentclass[10pt,a4paper]{amsart}
\usepackage[margin=19mm]{geometry}
\usepackage{amsmath,amssymb,mathtools}
\usepackage[scr=rsfs,cal=euler]{mathalfa}
\usepackage{xfrac}
\usepackage[unicode,hidelinks,bookmarks=false]{hyperref}
\newcommand{\ie}{i.e.\ }
\newcommand{\cf}{cf.\ }
\newcommand{\resp}{respectively\ }
\newcommand{\Subsection}{\subsection}
\providecommand{\nobreakdash}{\nobreak\hbox{-}}
\setlength{\emergencystretch}{2em}
\multlinegap0pt
\usepackage{bm}
\usepackage{bbm}
\usepackage{dsfont}
\usepackage{xspace}
\usepackage{cancel}
\usepackage{mathtools}
\usepackage{amsthm}
\makeatletter
\@ifundefined{theorem}{\newtheorem{theorem}{Theorem}}{}
\@ifundefined{lemma}{\newtheorem{lemma}[theorem]{Lemma}}{}
\makeatother
\newcommand{\PP}{\mathbb{P}}
\newcommand{\ind}{\mathds{1}}
\title[On the pointwise and sup-norm errors for local regression es]{On the pointwise and sup-norm errors for local regression estimators}
\author{J\'er\'emy Bettinger, Fran\c{c}ois Portier, Adrien Saumard}
\date{Source: arXiv:2507.07132v1 (CC BY 4.0)}
\begin{document}
\maketitle

\noindent\emph{Source and licence.} On the pointwise and sup-norm errors for local regression estimators. Version arXiv:2507.07132v1 (CC BY 4.0), distributed under the Creative Commons Attribution 4.0 International licence, as recorded in the arXiv licence metadata for this exact version and shown on the abstract page https://arxiv.org/abs/2507.07132v1. Full rights notice: licenses/arxiv-2507.07132-CC-BY-4.0.txt.\par\smallskip

\noindent\textit{Editorial scope.} Counted unit: the Lemma whose source label is \texttt{kappa} in the article (source0.tex lines 2159--2164, source label \texttt{kappa}), delivered together with the article's own complete original proof of it (source0.tex lines 2167--2210, one \texttt{proof} environment of 44 lines). No mathematics is written, repaired or re-derived: statement and proof are the article's own text line for line. Required context retained uncounted: none. Every definition, display and label of those lines is retained unchanged. Omitted from this package and not redistributed: 1--2158, 2165--2166, 2211--2452. Nothing inside a retained excerpt is omitted or altered: every retained body is delivered exactly as the article's lines are written, so no omission receipt of any kind accompanies this package. Citations: the retained excerpts carry no citation command at all, so no bibliography entry of the article is transcribed and no reference is redistributed. Reference adaptation: none was needed and none was made; every reference inside the delivered unit resolves inside it, so no unresolved reference or citation is produced. Printed slips kept literal: no printed word, symbol, hypothesis or number of the counted statement or of its proof was altered. Layout and class: the article's own class is \texttt{scrartcl} with packages including mathpazo, helvet, fontenc, babel, microtype, amsmath, amsfonts, amsthm, bm, dsfont, graphicx, enumitem, algorithm, algpseudocode; the wrapper is the lane's pinned \texttt{amsart} editorial wrapper at 10pt on A4, which restates the theorem-like environments the retained text needs and adds the article's own macro definitions that the retained text needs (\texttt{\textbackslash PP}, \texttt{\textbackslash ind}), with extra wrapper packages (bm, bbm, dsfont, xspace, cancel, mathtools). The delivered numbering is the unit's own. Assets: no retained span contains a figure environment, a graphics-inclusion command, a PDF-page inclusion, a table or a TikZ picture; no image file is copied into this package, the package declares no asset dependency and no image file is redistributed with it. Licence: version arXiv:2507.07132v1 of this article is distributed under the Creative Commons Attribution 4.0 International licence, as recorded in the arXiv licence metadata for this exact version and shown on the abstract page https://arxiv.org/abs/2507.07132v1. A full rights notice is delivered at licenses/arxiv-2507.07132-CC-BY-4.0.txt. Characters: every retained excerpt is pure ASCII, so the delivered engine, XeLaTeX with its default Latin Modern text font, reports no missing character.
\begin{lemma}\label{kappa}
     Let $W,W_1,\ldots, W_n$ be independent and identically distributed random variables on $\mathbb R_{+}$ with density $f_W$ and cumulative distribution function $F_W$. Suppose that there exists $c_0>0$ and $\kappa \geq 0$, such that  $ f_W (t) \leq c_0 F_W(t) ^{\kappa} $ for all $t\in  \mathbb R_{+} $. 
     Let $\delta \in (0,1) $. It holds, with probability at least $1-\delta $, 
   $$ W_{(2)} -W_{(1)} >    C^{-1} \delta  m^{ \kappa-1 }  ,$$ 
   where $C=  c_0 \Gamma(\kappa+1) 3^{\kappa+1}.$
\end{lemma}

\begin{proof}
    Note that 
\begin{align*}
 \PP (W_{(2)} -W_{(1)}   > t ) &= \sum_{k=1} ^ m \mathbb P (  W_j > W_k + t \, : \, \forall  j \neq k) \\
 &= \sum_{k=1} ^ m \mathbb E [  (1- F_W ( W_k + t ))^{m-1} ]\\
 &= m \mathbb E [   (1- F_W ( W + t ))^{m-1} ].
  \end{align*}
Let us define, 
%for a set $\mathcal{C}\subset \mathbb{R_+}$ and 
for any $t\in \mathbb{R}_+$,
\begin{equation*}
    D(t):= \frac{1}{m}\mathbb{P}(W_{(2)}-W_{(1)}>t) =  \mathbb E [   (1- F_W ( W + t ))^{m-1}  ].
   % \mathbb{1}_\mathcal{C}(W)
\end{equation*}
It holds $D(0)=1/m$, $D(+\infty)=0$ and by its definition through the deviation probability, $D$ is a non-increasing function on $\mathbb{R}_+$. 
%Assuming that $W$ has a density on $\mathbb{R}_+$ and that we can differentiate under the integral sign in the second expression of $D(t)$ - which is the case if for instance, the density $f_W$ of $W$ is bounded from above - then $D$ is differentiable on $\mathbb{R}_+$. More precisely, for any $t \in \mathbb{R}_+$,
%\begin{equation*}
 %   D^\prime(t)=-(m-1)\mathbb{E}[f_W(W+t)(1- F_W ( W + t ))^{m-2}].
  %  \end{equation*}
By using Fubini-Tonelli and integrating first with respect to $t$, we find
\begin{align*}
    &(m-1)\int_0^t \mathbb{E}[f_W(W+u)(1- F_W ( W + u ))^{m-2}] \, du \\ 
    &= \mathbb E\left[ (m-1)\int_0^t f_W(W+u)(1- F_W ( W + u ))^{m-2} \, du \right] \\ &= \mathbb E [   (1- F_W ( W ))^{m-1}  ] - \mathbb E [   (1- F_W ( W + t ))^{m-1}  ] \\ &= D(0)-D(t).
\end{align*}
We also have
\begin{align*}
      & \mathbb{E}[f_W(W+u)(1- F_W ( W + u ))^{m-2}] \\ &=   \int_0^{+\infty} f_W(r+u)f_W(r)(1- F_W ( r + u ))^{m-2} \, \ind_{f_W(r+u) > 0} \, dr \\
    & \leq  c_0 \int_0^{+\infty} f_W(r+u)  F_W(r)^\kappa (1- F_W ( r + u ))^{m-2} \, \ind_{f_W(r+u) > 0} \, dr.
\end{align*}
%Assume that $d\geq 2$. Note that the case $d=1$ is actually easier. 
Let us now apply a change of variable, which is justified because it is differentiable and bijective when $f_W$ is positive. Note also that $ F _W ( F_W^{-1} (u)) = u$ because $F_W$ is continuous. By setting $v=F_W(r+u)$, we get $dv=f_W(r+u)dr$, which gives, for any $\kappa \geq 0$,
\begin{align*}
    &\int_0^{+\infty} f_W(r+u)  F_W(r)^\kappa (1- F_W ( r + u ))^{m-2}\, \ind_{f_W(r+t) > 0} \, dr \\ &= \int_{F_W(u)}^{1}  F_W(F_W^{-1}(v)-u)^\kappa (1- v)^{m-2} \, \ind_{f_W(F^{-1}_W(v)) > 0} \, dv \\ &\leq  \int_{0}^{1} v^\kappa (1- v)^{m-2}dv 
\end{align*}
because $\ind_{f_W(F^{-1}_W(v)) > 0} \leq 1$ and $ F_W(F_W^{-1}(v)-u) \leq F_W(F_W^{-1}(v)) = v.$
Moreover,
\begin{align*}
\int_{0}^{1} v^\kappa (1- v)^{m-2}dv &\leq \int_{0}^{1} v^\kappa \exp(-v(m-2)) dv = \dfrac{1}{(m-2)^{\kappa+1}}\int_{0}^{m-2} s^\kappa e^{-s} ds \\ &\leq \dfrac{\Gamma(\kappa+1)}{(m-2)^{\kappa+1}}.
\end{align*}
Putting things together, we have shown that
$$D(0) - D(t) \leq c_0 \int_{0}^t (m-1)  \dfrac{\Gamma(\kappa+1)}{(m-2)^{\kappa+1}} du \leq m c_0 \Gamma(\kappa+1) \dfrac{3^{\kappa+1}}{m^{\kappa+1}} t = C \dfrac{t}{m^\kappa}.$$
In the latter upper-bound, we used $m-2 \geq m/3$, and $C = c_0 \Gamma(\kappa+1) 3^{\kappa+1}$. As a consequence, $$\PP\left(W_{(2)} -W_{(1)}   > t \right) = m D(t) \geq m D(0) - m \dfrac{C}{m^\kappa} t =   1 - \dfrac{C}{m^{\kappa-1}} t.$$
Choosing $t = { \delta  \, m^{\kappa-1}} / {C}$, leads to the statement.
\end{proof}

\end{document}
